\documentclass[12pt,letterpaper]{report}
\usepackage{graphicx}
\usepackage{scrextend}
\usepackage{vmargin}
\usepackage{graphicx}
\usepackage{multirow}
\usepackage[utf8]{inputenc}
\usepackage[spanish]{babel}
\usepackage{multicol}
\usepackage{enumerate}
\usepackage{float}
\usepackage{amsmath, amsthm, amssymb, amsfonts}
\usepackage[usenames]{color}
\parindent=0mm
\pagestyle{empty}
\definecolor{miorange}{rgb}{0.91, 0.43, 0.0}
\begin{document}
\setmargins{2.5cm}      
{1.5cm}                     
{2cm}  
{24cm}                    
{10pt}                          
{1cm}                          
{0pt}                             
{2cm}
\begin{flushright}
\today
\end{flushright}
\vspace{-1.75cm}
\section*{Tarea Clase 10}
\subsection*{Obtenga las soluciones de las siguientes ecuaciones:}
\begin{enumerate}
\item Calcular la altura de un edificio que proyecta una sombra de 6.5 m a la misma hora que un poste de 4.5 m de altura da una sombra de 0.90 m.
\item Los catetos de un triángulo rectángulo que miden 24 m y 10 m. ¿Cuánto medirán los catetos de un triángulo semejante al primero cuya hipotenusa mide 52 m?
\item Escribe de todas las formas posibles la ecuación de la recta que pasa por los puntos A(1,2) y B(-2,5)
\item Hallar la pendiente y la ordenada en el origen de la recta $3x+2y-7=0$.
\item La recta $r=3x+ny-7=0$ pasa por el punto A(3, 2) y es paralela a la recta $s=mx+2y-13=0$. Calcula m y n.
\end{enumerate}
\end{document}